\documentclass[10pt,a4paper]{amsart}
\usepackage[margin=19mm]{geometry}
\usepackage{amsmath,amssymb,mathtools}
\usepackage[scr=rsfs,cal=euler]{mathalfa}
\usepackage{xfrac}
\usepackage[unicode,hidelinks,bookmarks=false]{hyperref}
\newcommand{\ie}{i.e.\ }
\newcommand{\cf}{cf.\ }
\newcommand{\resp}{respectively\ }
\newcommand{\Subsection}{\subsection}
\providecommand{\nobreakdash}{\nobreak\hbox{-}}
\setlength{\emergencystretch}{2em}
\multlinegap0pt
\usepackage{caption}
\usepackage{dsfont}
\usepackage{mathtools}
\usepackage{multirow}
\usepackage{xcolor}
\newtheorem{remark}{Remark}
\newtheorem{lemma}{Lemma}
\usepackage{cleveref}
\crefname{remark}{Remark}{Remarks}
\crefname{lemma}{Lemma}{Lemmas}
\newcommand{\E}{\mathbf{E}}
\renewcommand{\P}{\mathbf{P}}
\newcommand\wh[1]{\widehat{#1}}
\title{On the principal eigenvectors of random Markov matrices}
\author{Jacob Calvert, Frank den Hollander,, Dana Randall}
\date{Source: arXiv:2505.01608v2 (CC BY 4.0)}
\begin{document}
\maketitle

\noindent\emph{Source and licence.} On the principal eigenvectors of random Markov matrices. Version arXiv:2505.01608v2 (CC BY 4.0), distributed by arXiv under the Creative Commons Attribution 4.0 International licence, http://creativecommons.org/licenses/by/4.0/, as recorded in the arXiv licence metadata for this exact version and shown on the abstract page https://arxiv.org/abs/2505.01608v2. Full licence notice: \texttt{licenses/arxiv-2505.01608-CC-BY-4.0.txt}.\par\smallskip

\noindent\textit{Editorial scope.} Counted unit: the article's lemma l\_infty (source0.tex lines 284--289). The article's complete original proof of it is delivered with it (source0.tex lines 293--348, one \texttt{proof} environment of 56 lines, which the article places away from the statement). No mathematics is written, repaired or re-derived: the statement and the proof are the article's own lines, in the article's own notation, and no retained line is substituted or re-flowed.  Retained uncounted as the source-local context the proof needs: lines 122--124; lines 200--205; lines 234--239; lines 259--264.  Nothing was omitted from any retained span, so the package carries no omission record.  No retained line contains a citation, so the wrapper carries no bibliography and no bibliography entry.  No retained line contains a figure environment, an image-inclusion command or drawing code, so no image is delivered and the package declares no asset dependency.  Editorial formatting only, in addition: an AMS \texttt{amsart} preamble that restates the article's own declarations and macros used by the retained lines, namely the environments \texttt{remark}, \texttt{lemma} on separate counters, together with the standard packages those lines need.  The retained environments print in this document as Remark 1, Lemma 1 in order of appearance, whereas the article prints the same items with the numbering of its own full document.  The complete native source of the article as distributed in the arXiv e-print for this exact version is bundled unchanged as \texttt{sources/177515/source0.tex}.
\begin{remark}\label{rem: A prim}
    For the rest of the paper, we will work on the event that $A$ and $\wh{A}$ are primitive, which holds for all sufficiently large $n$ a.s. On this event, the Markov matrices $Q$, $P$, and $\wh{Q}$ are well-defined and irreducible, and hence the corresponding invariant distributions $\pi_Q$, $\pi_P$, and $\pi_{\wh Q}$ exist and are unique. Additionally, the kernels $P$ and $\wh{Q}$ are aperiodic.
\end{remark}

\begin{lemma}\label{lem:bai yin}
Let $(X_{i,j})_{i, j \geq 1}$ be an array of i.i.d.\ random variables. For any real numbers $\alpha \in (1/2,1]$, $\beta \geq 0$, and $C > 0$, if $\E(|X_{1,1}|^{(1+\beta)/\alpha}) < \infty$, then
\[
		\max_{i \in [Cn^\beta]} \left| \sum_{j \in [n]} (X_{i,j} - \E (X_{1,1})) \right| = o(n^{\alpha}) \quad \mathrm{a.s.}
\]
\end{lemma}

\begin{lemma}\label{lem:l2 control}
Let $K$ denote either $\wh{Q}$ or $P$. If $\mathcal{L}$ has finite fourth moment, then
\[
	\max_{i \in [n]} \sum_{j \in [n]} K_{i,j}^{\,2} = O(n^{-1}) \quad \mathrm{a.s.}
\]
\end{lemma}

\begin{lemma}\label{lem:max control}
Let $K$ denote $\wh{Q}$ or $P$. If $\mathcal{L}$ has finite $p$-th moment for some $p >4$, then, for every $a < 1 - \frac{2}{p}$,
\[
	\max_{i,j \in [n]} K_{i,j} = O (n^{-a}) \quad \mathrm{a.s.}
\]
\end{lemma}

\begin{lemma}\label{lemma:P l_infty}
If $\mathcal{L}$ has finite $p$-th moment for some $p > 4$, then, for every $a < \min\{\frac12,1-\frac{4}{p}\}$,
\begin{equation}\label{eq:P l_infty}
	\| \pi_{\wh{Q}} - u \|_\infty = O(n^{-(1+a)}) \quad \mathrm{a.s.}
\end{equation}
\end{lemma}

\begin{proof}[Proof of \cref{lemma:P l_infty}]
	Fix $p > 4$ and $a < \min\{\frac12,1-\frac{4}{p}\}$. According to \cref{rem: A prim}, $\wh{Q}$ is irreducible and aperiodic, and hence $\pi_{\wh{Q}} = \pi_{\wh{Q}^2}$ for all sufficiently large $n$ a.s. It therefore suffices to show that
\begin{equation}\label{eq:p2ik lb}
	\max_{i,k \in [n]} \left| (\wh{Q}^2)_{i,k} - \frac1n \right| = O\big(n^{-(1+a)} \big) \quad \mathrm{a.s.}
\end{equation}
Indeed, because $\pi_{\wh{Q}}$ is stationary, on the event in \cref{eq:p2ik lb}, it satisfies
\[
	\pi_{\wh{Q}} (i) = \sum_{j \in [n]} \pi_{\wh{Q}} (j) (\wh{Q}^2)_{j,i} = \frac{1}{n} + O(n^{-(1+a)}),
\]
which implies \cref{eq:P l_infty}.

To prove \cref{eq:p2ik lb}, we begin by noting that, conditionally on the $i$-th row $\wh{Q}_{i,\cdot}$ of $\wh{Q}$, the random variables $Y_j = \wh{Q}_{i,j} (\wh{Q}_{j,k} - \frac{1}{n-1})$ indexed by $j \in [n]\setminus\{i,k\}$ are independent, centered, and satisfy
\begin{equation}\label{eq:sumy}
	S_n (i,k) = \sum_{j \in [n]\setminus\{i,k\}} Y_j = (\wh{Q}^2)_{i,k} - \frac{1-\wh{Q}_{i,k}}{n-1}. 
\end{equation}
Bernstein's inequality states that there is a universal constant $C > 0$ such that, for every $t > 0$,
\begin{equation}\label{eq:bernstein}
	\P \left( |S_n (i,k)| > t \mid \wh{Q}_{i,\cdot} \right) \leq 2 \exp \left(- C \min\left\{\frac{t^2}{v_n}, \frac{t}{b_n} \right\} \right),
\end{equation}
with
\[
	v_n = \sum_{j \in [n]\setminus\{i,k\}} \E (Y_j^2 \mid \wh{Q}_{i,\cdot}), \qquad b_n = \max_{j \in [n]\setminus\{i,k\}} |Y_j|.
\]
Regarding $v_n$, \cref{lem:bai yin} implies that, because $\mathcal{L}$ has finite second moment, the off-diagonal row sums $R_i = \sum_{j:\, j \neq i} X_{i,j}$ satisfy
\[
	\min_{i \in [n]} R_i = (1-o(1))\, \E(X_{1,1}) n \quad \mathrm{a.s.}
\]
On this event, as well as the a.s.\ event of \cref{lem:l2 control}, the fact that $X_{j,k}$ is independent of $\wh{Q}_{i,\cdot}$ and has finite second moment implies that
\[
	v_n = \sum_{j \in [n]\setminus\{i,k\}} \wh{Q}^2_{i,j} \, \E (\wh{Q}^2_{j,k} \mid \wh{Q}_{i,\cdot}) = O(n^{-2}) \sum_{j \in [n]\setminus\{i,k\}} \wh{Q}^2_{i,j} = O(n^{-3}).
\]
Furthermore, on the a.s.\ event of \cref{lem:max control}, for every $s < 2(1 - \frac{2}{p})$, 
\[
	b_n \leq \max_{j \in [n]\setminus\{i,k\}} \wh{Q}_{i,j} \Big(\wh{Q}_{j,k} + \frac{1}{n-1}\Big) = O(n^{-s}).
\]
In particular, we choose $s > 1+a$.

When $t_n = \Theta (n^{-(1+a)})$ the rate in \cref{eq:bernstein} satisfies
\[
	\min\left\{\frac{t_n^2}{v_n}, \frac{t_n}{b_n} \right\} = \Omega \left(\min \left\{ n^{1-2a}, n^{s-1-a}\right\}\right) = \Omega (n^{\delta}) \quad \mathrm{a.s.},
\]
where $\delta = \min\{1-2a,s-1-a\}> 0$. A union bound over the $n^2$ possible values of $i$ and $k$, together with \cref{eq:bernstein}, imply that
\[
	\P \left( \max_{i,k \in [n]} |S_n (i,k)| > t_n \right) \leq 2n^2 \exp \big(- \Omega(n^\delta) \big) \quad \mathrm{a.s.}
\]
Because this upper bound is summable, the Borel--Cantelli lemma implies that
\[
	\max_{i,k \in [n]} |S_n (i,k)| = O (n^{-(1+a)}) \quad \mathrm{a.s.}
\]

To relate this to \cref{eq:p2ik lb}, we combine the preceding fact with \cref{eq:sumy}. The triangle inequality and \cref{lem:max control} imply that, for every $c < 2(1-\frac{1}{p})$,
\[
	\max_{i,k \in [n]} \left| (\wh{Q}^2)_{i,k} - \frac1n \right| \leq \max_{i,k \in [n]} \left( \left| S_n (i,k) \right| + \left| \frac{1-\wh{Q}_{i,k}}{n-1} - \frac1n \right| \right) = O (n^{-(1+a)}) + O(n^{-c}) \quad \mathrm{a.s.},
\]
In particular, we can take $c = \frac{3}{2} > 1 + a$ because $p > 4$, in which case $O(n^{-c}) = o(n^{-(1+a)})$. This proves \cref{eq:p2ik lb}, and completes the proof.
\end{proof}

\end{document}
